
\begin{frame}[fragile]{Retour en arrière}

Notre requête sur les films avec John Ferguson, après optimisation algébrique.

\vfill
\begin{small}
\psset{levelsep=1cm,treesep=0.5cm}
\IteratorSansNom{$\pi_{titre}$}{}{
     \IteratorSansNom{$\Join$}{}{ 
       \IteratorSansNom{$\sigma_{annee=1958}$}{}{
       \Database{Film}{}
       }
  \IteratorSansNom{$\sigma_{nom\_role='Ferguson'}$}{}{
       \Database{Role}{}
    }
   }
}
\end{small}


\vfill

Il nous restait à choisir les chemins d'accès et les algorithmes de jointure.


    
\end{frame}

\begin{frame}{Continuons à optimiser notre requête}

Hyp.: le système 
cherche à utiliser les index pour les jointures .

\vfill

\[
  Film \Join_{id=id\_film} Role
\]
  
ou 

\[
   Role \Join_{id\_film=id} Film
\]

\vfill

Après analyse des index, la bonne forme est
\[
   \pi_{titre}(\sigma_{nom\_role='John\ Ferguson'}(Role)  \Join_{id\_film=id}  \sigma_{annee=1958} (Film) )
\]

\end{frame}

\begin{frame}{Le plan d'exécution}

\figSlideWithSize{planEx-full1}{9cm}

\end{frame}


\begin{frame}{Et si je n'ai pas d'index?}

\figSlideWithSize{planEx-full-hash}{7cm}

\end{frame}
\begin{frame}{Résumé: la phase d'optimisation}


\vfill
\figSlideWithSize{exec-optim-4}{14cm}
\vfill


\only<2|handout:2>{
\centerline{\alert{Merci!}}
}

\end{frame}


\begin{comment}
\begin{frame}{Et si j'indexe le nom du rôle?}

\figSlideWithSize{planEx-full2}{5.5cm}

\end{frame}

\begin{frame}{Et si j'indexe la clé étrangère?}

\figSlideWithSize{planEx-full3}{5.5cm}

\end{frame}
\end{comment}